\section{Separate first-exit minimality from tile geometry}
\label{sec:firstexit}

\subsection{Two proofs with the same well-ordering step}

\code{WangTiles942InitialStrip.left_wall_exists} finds the first tile to
the left of Start that is not the left initializer. Its right-hand analogue
finds the first tile after the head that is not the right initializer.
Both construct a \code{Nat.find} witness, prove a bound from a known
failure, prove that every preceding index is still an initializer, and
split between an immediate exit and an initializer predecessor
\cite{strip}.

The geometric data are different and should stay different. The left
search starts at one and has a witness by \(N\); the right search starts at
two and has a witness by \(N+1\). The left seam is classified by east color
two, while the right seam is classified by west color four. The proposed
shared theorem captures only the well-ordering argument.

\subsection{The common theorem}

For a decidable predicate \(P:\N\to\mathrm{Prop}\), a starting index
\(a\), and a known failing index \(b\ge a\), prove
\[
\begin{split}
\exists k,\quad &a\le k\le b,\quad \neg P(k),\\
&\forall j,\ a\le j<k\ \Longrightarrow\ P(j),\\
&k=a\quad\lor\quad P(k-1).
\end{split}
\]
The exact interface is:
\begin{lean}
theorem exists_first_failure (P : Nat → Prop) [DecidablePred P]
    (a b : Nat) (hab : a ≤ b) (hb : ¬ P b) :
    ∃ k : Nat, a ≤ k ∧ k ≤ b ∧ ¬ P k ∧
      (∀ j : Nat, a ≤ j → j < k → P j) ∧
      (k = a ∨ P (k - 1))
\end{lean}

The disjunction is necessary. If the first position already fails,
\(k=a\), then \(k-1\) need not lie in the searched interval. When
\(a=0\), natural subtraction truncates, making an unconditional
predecessor conclusion particularly inappropriate. The focused checks
include this zero-start immediate-exit boundary and a later-exit case.

The proof chooses the least index satisfying \(a\le j\land\neg P(j)\).
Minimality yields the interior condition. If the chosen index differs from
\(a\), then \(a\le k-1<k\), so the same interior condition proves the
predecessor case. Periodicity, finite cardinality, and tiles are absent
from this theorem.

\subsection{Applications retain the exact geometric premises}

For the left wall, instantiate
\[
P(j):=T(-j)=\mathrm{initLeft},\qquad a=1,\quad b=N.
\]
The original periodic-Start argument supplies \(\neg P(N)\). For the
right wall, instantiate
\[
P(j):=T(j)=\mathrm{initRight},\qquad a=2,\quad b=N+1.
\]
The original periodic-head argument supplies \(\neg P(N+1)\).

The \code{k = a} branch uses Start or the head initializer respectively.
The other branch uses the returned initializer predecessor. Actual
horizontal matching then supplies the boundary edge color, and the
existing local classification theorem identifies the seam. No uniqueness
of Start, global rectangular tiling, or machine-run hypothesis is added.

\subsection{An explicit source-size tradeoff}

The client module decreases from 148 to 133 physical lines. The new
production helper is 28 lines, so the aggregate changes from 148 to 161:
\emph{13 additional physical lines}. This is not a net shortening claim.
The gain is that two copies of the least-exit proof become one theorem,
and each geometric proof contains only its own boundary reasoning.

For a project optimizing immediate total source length, defer this patch.
For a project expecting further initial-strip or reset-strip arguments,
the shared invariant can be worthwhile. The archive keeps this patch
separate, so that decision does not affect the other reductions.

This is distinct from the older finite-invariant predecessor proposal:
that proposal derives backward preservation from an injective map of a
finite invariant set. Here the shared mechanism is least-failure
minimality on a bounded natural interval, with an explicit starting
boundary. Neither finite injectivity nor an orbit-period argument enters
the helper.
